\documentclass[10pt]{article}
\usepackage[T1]{fontenc}
\usepackage[utf8]{inputenc}
\usepackage[french]{babel}
\usepackage{amsmath,amssymb,mathrsfs,mathpazo}
\usepackage{geometry}
\geometry{a4paper,margin=2cm}
\usepackage{graphicx,array,tabularx,tabularray,colortbl,multirow,multicol,booktabs,enumitem,pifont,ulem,fancybox,diagbox,cancel}
\usepackage{pstricks,pstricks-add,pst-plot,pst-tree,pst-node,pst-text,pst-func,pst-math,pst-bspline,pst-3dplot,pst-eucl,pst-solides3d}
\let\pscalculate\undefined
\usepackage{pst-bezier}
\usepackage{tikz,pgfplots,tkz-tab}
\usetikzlibrary{arrows,arrows.meta,patterns,calc,positioning,trees,decorations.markings,shapes.geometric,shapes.misc,shapes.arrows,intersections,decorations.pathmorphing,decorations.pathreplacing,backgrounds,shadows,babel}
\pgfplotsset{compat=1.18}
\usepgfplotslibrary{fillbetween}
\usepackage[np]{numprint}
\usepackage{listings,makecell,pas-tableur,twoopt}
\let\barre\undefined
\usepackage[pstricks]{tablvar}
\definecolor{darkgreen}{RGB}{0,100,0}
\providecommand{\red}{\color{red}}
\providecommand{\blue}{\color{blue}}
\providecommand{\green}{\color{green}}
\providecommand{\black}{\color{black}}
\providecommand{\white}{\color{white}}
\providecommand{\gray}{\color{gray}}
\providecommand{\cyan}{\color{cyan}}
\providecommand{\magenta}{\color{magenta}}
\providecommand{\yellow}{\color{yellow}}
\providecommand{\violet}{\color{violet}}
\providecommand{\orange}{\color{orange}}
\providecommand{\euro}{\texteuro}
\providecommand{\EUR}{\texteuro}
\providecommand{\ieme}{\textsuperscript{e}}
\providecommand{\degres}{\ensuremath{^\circ}}
\providecommand{\degree}{\ensuremath{^\circ}}
\providecommand{\vv}[1]{\overrightarrow{#1}}
\providecommand{\Coord}[1]{#1}
\providecommand{\decosix}{}
\providecommand{\starredbullet}{\ensuremath{\bullet}}
\pagestyle{empty}
\setlength{\parindent}{0pt}
\usepackage{adjustbox}\DeclareUnicodeCharacter{2236}{:}\DeclareUnicodeCharacter{2190}{\ensuremath{\leftarrow}}
\providecommand{\R}{\mathbb{R}}
\providecommand{\N}{\mathbb{N}}
\providecommand{\Z}{\mathbb{Z}}
\providecommand{\Q}{\mathbb{Q}}
\providecommand{\D}{\mathbb{D}}
\providecommand{\C}{\mathbb{C}}
\providecommand{\e}{{\mathrm{e}}}
\providecommand{\vect}[1]{\overrightarrow{#1}}
\providecommand{\vectt}[1]{\overrightarrow{\mathrm{#1}}}
\providecommand{\vvt}[1]{\overrightarrow{\mathrm{#1}}}
\providecommand{\np}[2][]{{#2~#1}}
\providecommand{\vv}[1]{{\overrightarrow{#1}}}
\renewcommand{\d}{{\mathrm{d}}}
\providecommand{\degres}{{^\circ}}
\providecommand{\degre}{{^\circ}}
\providecommand{\blue}{\color{blue}}
\providecommand{\red}{\color{red}}
\providecommand{\enskip}{\quad}
\providecommand{\llbracket}{\lbrack\!\lbrack}
\providecommand{\rrbracket}{\rbrack\!\rbrack}
\providecommand{\euro}{{\text{€}}}
\providecommand{\no}{{\text{n°}}}
\providecommand{\interval}[2]{{[#1;#2]}}
\providecommand{\intervFO}[2]{{[#1;#2[}}
\providecommand{\dotfill}{\ldots}

\providecommand{\ds}{\displaystyle}
\providecommand{\textding}[1]{#1}
\providecommand{\Oij}{\(\mathrm{O};\vec{\imath},\vec{\jmath}\)}
\providecommand{\Oijk}{\(\mathrm{O};\vec{\imath},\vec{\jmath},\vec{k}\)}
\providecommand{\Ouv}{\(\mathrm{O};\vec{u},\vec{v}\)}
\providecommand{\intervFF}[2]{\ensuremath{[#1;#2]}}
\renewcommand{\euro}{\ifmmode\text{\texteuro}\else\texteuro\fi}
\DeclareUnicodeCharacter{00B0}{\ifmmode^{\circ}\else\textdegree\fi}
\DeclareUnicodeCharacter{20AC}{\ifmmode\text{\texteuro}\else\texteuro\fi}
% Énoncés : Éducation nationale. Transcription et corrigés : APMEP, Denis Vergès et contributeurs. Découpage et classement : CoopMaths, Sébastien Lozano, Éric Elter et contributeurs. Adaptation de la présentation : Maths Explorer.
% CC BY-SA : https://coopmaths.fr/www/about/
% https://coopmaths.fr/alea/static/bac/2022/tex/bac_2022_09_sujet2_ameriquesud_2.tex
\begin{document}
\noindent\textit{APMEP / CoopMaths. CC BY-SA. Présentation adaptée par Maths Explorer.}\par\medskip
\medskip

Le but de cet exercice est d'étudier la fonction $f$, définie sur $]0~;~ +\infty[$ ,par:

\[f(x) =3x - x \ln (x) - 2 \ln (x).\]

\smallskip

\textbf{PARTIE A: Étude d'une fonction auxiliaire }\boldmath $g$ \unboldmath

\medskip

Soit $g$ la fonction définie sur $]0~;~ +\infty[$ par 

\[g(x) = 2(x - 1) - x \ln (x).\]

On note $g'$ la fonction dérivée de $g$. On admet que $\displaystyle\lim_{x \to + \infty} g(x) = - \infty$.

\medskip

\begin{enumerate}
\item Calculer $g(1)$ et $g(\text{e})$.
\item Déterminer $\displaystyle\lim_{x \to + 0} g(x)$ en justifiant votre démarche.
\item Montrer que, pour tout $x > 0$,\: $g'(x) = 1 - \ln (x)$.

En déduire le tableau des variations de $g$ sur $]0~;~ +\infty[$.
\item Montrer que l'équation $g(x) = 0$ admet exactement deux solutions distinctes sur
$]0~;~ +\infty[$ : 1 et $\alpha$ avec $\alpha$ appartenant à l'intervalle $[\text{e}~;~+\infty[$.

On donnera un encadrement de $\alpha$ à 0,01 près.
\item En déduire le tableau de signes de $g$ sur $]0~;~ +\infty[$. 
\end{enumerate}

\bigskip

\textbf{PARTIE B : Étude de la fonction} \boldmath $f$ \unboldmath

\medskip

On considère dans cette partie la fonction $f$, définie sur $]0~;~ +\infty[$,par 

\[f(x) = 3x - x \ln (x)- 2\ln (x).\]

On note $f'$ la fonction dérivée de $f$.

La représentation graphique $\mathcal{C}_f$ de cette fonction $f$ est donnée dans le repère \Oij{} ci-dessous. On admet que : $\displaystyle\lim_{x \to 0}f(x) = + \infty$.

\begin{center}
\psset{unit=0.82cm,arrowsize=2pt 3}
\begin{pspicture*}(-1,-2.01)(16,6)
\psgrid[gridlabels=0pt,subgriddiv=1,gridwidth=0.15pt](0,-2)(16,6)
\psaxes[linewidth=1.25pt,labelFontSize=\scriptstyle]{->}(0,0)(0,-1.95)(16,5.95)
\psaxes[linewidth=1.25pt,labelFontSize=\scriptstyle](0,0)(0,-1.95)(16,5.95)
\psplot[plotpoints=2000,linewidth=1.25pt,linecolor=blue]{0.01}{16}{x 3 mul x ln x 2 add mul sub}
\uput[ur](10.5,2){\blue $\mathcal{C}_f$}
\uput[u](15.75,0){$x$}
\uput[l](0,5.75){$y$}
\end{pspicture*}
\end{center}

\begin{enumerate}
\item Déterminer la limite de $f$ en $+ \infty$ en justifiant votre démarche.
\item 
	\begin{enumerate}
		\item Justifier que pour tout $x > 0$,\: $f'(x) = \dfrac{g(x)}{x}$.
		\item En déduire le tableau des variations de $f$ sur $]0~;~ +\infty[$.
	\end{enumerate}	
\item On admet que, pour tout $x > 0$, la dérivée seconde de $f$, notée $f''$, est définie par $f''(x) = \dfrac{2 - x}{x^2}$.

Étudier la convexité de $f$ et préciser les coordonnées du point
 d'inflexion de $\mathcal{C}_f$.
\end{enumerate}

\bigskip

\end{document}
